\section{Proxy-to-mechanism validation}
\label{sec:methodological-discipline}

One principle recurs across every intervention above, and is applied to the
report's own claims. For each intervention two questions are asked
separately: (i) \emph{did it change its intended proxy?} and (ii) \emph{did
the anatomical mechanism the proxy stands for actually improve?} The
investigation repeatedly found $\text{proxy}\uparrow \not\Rightarrow
\text{mechanism}\uparrow$: lower Chamfer with a worse skeleton (\F{F1}), a
scale statistic that improved without the anterior morphology improving,
and a validation instrument whose own synthetic estimate did not survive
expert ground truth (\S\ref{sec:reliability}). This is a general lesson of
optimisation and scientific machine learning: an objective is a proxy for
the property one cares about, and is treated here as a design rule:
every result is reported with the level of evidence that supports it
(Table~\ref{tab:evidence}).
